
\chapter*{参考文献}
\phantomsection
\addcontentsline{toc}{chapter}{参考文献}

\begingroup
\def\chapter*#1{}%
\newcommand{\sameauthor}{\rule[-1pt]{1.5em}{0.4pt}}%
\begin{thebibliography}{[MO129242]}
\sloppy\Urlmuskip=0mu plus 1mu\relax%
\hbadness=99999 \vbadness=99999 \hfuzz=20pt \vfuzz=5pt%

\bibitem[Al]{bib-Al}
P. Aluffi, \emph{Algebra, Chapter 0}, Grad. Stud. in Math. \textbf{104}, Amer.
Math. Soc., Providence, RI, 2009.

\bibitem[Ar1]{bib-Ar1}
M. Artin, ``On the solutions of analytic equations,'' Invent. Math. \textbf{5}
(1968), no. \textbf{4}, 277--291.

\bibitem[Ar2]{bib-Ar2}
\sameauthor, ``Algebraic approximation of structures over complete local
rings,'' Publ. Math. IHES \textbf{36} (1969), 23--58.

\bibitem[Ar3]{bib-Ar3}
\sameauthor, \emph{Algebra}, Prentice Hall, Inc., Englewood Cliffs, NJ, 1991.

\bibitem[At1]{bib-At1}
M. F. Atiyah, ``On analytic surfaces with double points,'' Proc. Roy. Soc.
London Ser. A \textbf{247} (1958), 237--244.\footnote{原书此处作『237—244』
（长横），当为『237--244』之误，径改。——译者注}

\bibitem[At2]{bib-At2}
\sameauthor, ``Mathematics in the 20th century,'' Amer. Math. Monthly
\textbf{108} (2001), no. \textbf{7}, 654--666.

\bibitem[AtM]{bib-AtM}
M. F. Atiyah and I. G. Macdonald, \emph{Introduction to Commutative Algebra},
Addison-Wesley Publ. Co., Reading, MA, 1969.

\bibitem[Ba]{bib-Ba}
L. Bădescu, \emph{Algebraic Surfaces}, V. Maşek trans., Universitext,
Springer-Verlag, New York, 2001.

\bibitem[BCDKT]{bib-BCDKT}
M. Baker, B. Conrad, S. Dasgupta, K. Kedlaya, and J. Teitelbaum, \emph{p-adic
geometry}, D. Savitt and D. Thakur ed., Univ. Lect. Series \textbf{45}, Amer.
Math. Soc., Providence, RI, 2008.

\bibitem[BPV]{bib-BPV}
W. Barth, C. Peters, and A. van de Ven, \emph{Compact Complex Surfaces},
Ergeb. Math. Grenz. \textbf{4}, Springer-Verlag, Berlin, 1984.

\bibitem[Be]{bib-Be}
A. Beauville, \emph{Complex Algebraic Surfaces}, 2nd ed., R. Barlow trans.,
Cambridge U. P., Cambridge, 1996.

\bibitem[BL]{bib-BL}
C. Birkenhake and H. Lange, \emph{Complex Abelian Varieties}, 2nd ed., Grund.
Math. Wiss. \textbf{302}, Springer-Verlag, Berlin, 2004.

\bibitem[BLR]{bib-BLR}
S. Bosch, W. Lütkebohmert, and M. Raynaud, \emph{Néron Models}, Ergeb. Math.
Grenz. \textbf{21}, Springer-Verlag, Berlin, 1990.

\bibitem[Bo]{bib-Bo}
N. Bourbaki, \emph{Éléments de mathématique, Fasc. XXX: Algèbre Commutative},
Actualités Scientifiques et Industrielles, no. \textbf{1308}, Hermann, Paris,
1964.

\bibitem[BGS]{bib-BGS}
S. Boissière, O. Gabber, and O. Serman, ``Sur le produit de variétés
localement factorielles ou $\mathbb{Q}$-factorielles,'' arxiv:1104.1861v1,
preprint 2011.

\bibitem[BroP]{bib-BroP}
R. Brown and T. Porter, ``Analogy, concepts and methodology, in mathematics,''
UWB Math Preprint, May 26, 2006.

\bibitem[BruH]{bib-BruH}
W. Bruns and J. Herzog, \emph{Cohen-Macaulay Rings}, Cambridge U. P.,
Cambridge, 1993.

\bibitem[CL]{bib-CL}
F. Call and G. Lyubeznik, ``A simple proof of Grothendieck's theorem on the
parafactoriality of local rings,'' in \emph{Commutative algebra: syzygies,
multiplicities, and birational algebra (South Hadley, MA, 1992)}, Contemp.
Math. \textbf{159}, Amer. Math. Soc., Providence, RI, 1994, 15--18.

\bibitem[Carr]{bib-Carr}
L. Carroll, \emph{Alice's Adventures in Wonderland}, Macmillan, London, 1865.

\bibitem[ChJLO]{bib-ChJLO}
R. Cheng, L. Ji, M. Larson, and N. Olander, ``Theorem of the base,'' in
\emph{Stacks Project Expository Collection (SPEC)}, 163--193, London Math.
Soc. Lecture Note Ser. \textbf{480}, Cambridge U. P., Cambridge, 2022.

\bibitem[Coh]{bib-Coh}
G. Cohen, \emph{Essai d'explication du Cimetière Marin}, Librairie Gallimard,
1933.

\bibitem[Con]{bib-Con}
B. Conrad, ``Deligne's notes on Nagata compactifications,'' J. Ramanujan Math.
Soc. \textbf{22} (2007), no. \textbf{3}, 205--257.

\bibitem[CP]{bib-CP}
V. Cossart and O. Piltant, ``Resolution of singularities of arithmetical
threefolds,'' J. Algebra \textbf{529} (2019), no. \textbf{7}, 268--535.

\bibitem[Dan]{bib-Dan}
V. I. Danilov, ``Samuel's conjecture,'' Mat. Sb. (N.S.) \textbf{81}
(\textbf{123}) (1970), 132--144.

\bibitem[Dav]{bib-Dav}
E. D. Davis, ``Ideals of the principal class, R-sequences and a certain
monoidal transformation,'' Pacific J. Math. \textbf{20} (1967), 197--205.

\bibitem[DW]{bib-DW}
R. Dedekind and H. Weber, \emph{Theory of algebraic functions of one
variable}, J. Stillwell trans. History of Math. \textbf{39}, Amer. Math.
Soc., Providence, RI; London Math. Soc. London, 2012.

\bibitem[De]{bib-De}
P. Deligne, ``Courbes elliptiques: formulaire d'après Tate,'' in \emph{Modular
Functions of One Variable IV} (Proc. Internat. Summer School, Univ. Antwerp,
Antwerp, 1972), 53--73, Lect. Notes in Math. \textbf{476}, Springer-Verlag,
Berlin, 1975.

\bibitem[DeI]{bib-DeI}
P. Deligne and L. Illusie, ``Relèvements modulo $p^2$ et décomposition du
complexe de de Rham,'' Invent. Math. \textbf{89} (1987), no. \textbf{2},
247--270.

\bibitem[Di]{bib-Di}
J. Dieudonné, ``On regular sequences,'' Nagoya Math. J. \textbf{27} (1966),
355--356.

\bibitem[Do]{bib-Do}
I. Dolgachev, ``Luigi Cremona and cubic surfaces,'' in \emph{Luigi Cremona
(1830--1903)}, 55--70, Incontr. Studio, \textbf{36}, Instituto Lombardo di
Scienze e Lettere, Milan, 2005.

\bibitem[DS]{bib-DS}
R. Donagi and R. C. Smith, ``The structure of the Prym map,'' Acta Math.
\textbf{146} (1981), no. \textbf{1--2}, 25--102.

\bibitem[DZ]{bib-DZ}
H. Duan and X. Zhao, ``Schubert calculus and the integral cohomology of
exceptional Lie groups,'' arXiv:0711.2541v10, preprint 2007.

\bibitem[DF]{bib-DF}
D. Dummit and R. Foote, \emph{Abstract Algebra}, 3rd ed., Wiley and Sons,
Inc., Hoboken, NJ, 2004.

\bibitem[E]{bib-E}
D. Eisenbud, \emph{Commutative Algebra with a View to Algebraic Geometry},
Grad. Texts in Math. \textbf{150}, Springer-Verlag, New York, 1995.

\bibitem[EGH]{bib-EGH}
D. Eisenbud, M. Green, and J. Harris, ``Cayley-Bacharach theorems and
conjectures,'' Bull. AMS \textbf{33} (1996), no. \textbf{3}, 295--324.

\bibitem[Ek]{bib-Ek}
P. Eklof, ``Lefschetz's principle and local functors,'' Proc. Amer. Math.
Soc. \textbf{37} (1973), no. \textbf{2}, 1973.\footnote{疑当为『197--204』
（页码），兹照录原书。又，期号原书作 no.~2，通行著录为 no.~1，亦疑有误，兹照录。——译者注}

\bibitem[EE]{bib-EE}
E. Enochs and S. Estrada, ``Relative homological algebra in the category of
quasi-coherent sheaves,'' Adv. Math. \textbf{194} (2005), 284--295.

\bibitem[FGIKNV]{bib-FGIKNV}
B. Fantechi, L. Göttsche, L. Illusie, S. Kleiman, N. Nitsure and A. Vistoli,
\emph{Fundamental algebraic geometry: Grothendieck's FGA explained},
Mathematical Surveys and Monographs \textbf{123}, Amer. Math. Soc.,
Providence, RI, 2005.

\bibitem[FK]{bib-FK}
H. M. Farkas and I. Kra, \emph{Riemann Surfaces}, Grad. Texts in Math.
\textbf{71}, Springer-Verlag, New York, 1980.

\bibitem[FLS]{bib-FLS}
R. Feynman, R. Leighton, and M. Sands, \emph{The Feynman Lectures on Physics,
Volume III}, Addison Wesley, 1971.

\bibitem[Fi]{bib-Fi}
F. Scott Fitzgerald, ``The Crack-Up: A desolately frank document from one for
whom the salt of life has lost its savor,'' Esquire Magazine, February 1936,
Esquire Inc., Chicago, IL.

\bibitem[FR]{bib-FR}
G. Frey and H.-G. Rück, ``The strong Lefschetz principle in algebraic
geometry,'' manuscripta math. \textbf{55} (1986), no. \textbf{3--4},
385--401.

\bibitem[F1]{bib-F1}
W. Fulton, \emph{Young Tableaux: With Applications to Representation Theory
and Geometry}, Cambridge U. P., Cambridge, 1997.

\bibitem[F2]{bib-F2}
\sameauthor, \emph{Intersection Theory}, 2nd ed., Ergeb. Math. Grenz.
\textbf{2}, Springer-Verlag, New York, 1998.

\bibitem[Ge]{bib-Ge}
S. Germain, \emph{Œuvres Philosophiques de Sophie Germain} (nov. éd.),
Librairie de Firmin-Didot et Cie, Paris, 1896.\footnote{原书此处作『Librarie』，
当为『Librairie』之误，径改。——译者注}

\bibitem[GS]{bib-GS}
P. Gille and T. Szamuely, \emph{Central Simple Algebras and Galois
Cohomology}, Cambridge U. P., Cambridge, 2006.

\bibitem[GW]{bib-GW}
U. Görtz and T. Wedhorn, \emph{Algebraic Geometry I}, Vieweg + Teubner,
Wiesbaden, 2010. (Appendices C, D, and E, collecting properties of schemes and
morphisms and their relationships, are alone worth the price.)

\bibitem[GR]{bib-GR}
H. Grauert and R. Remmert, \emph{Coherent Analytic Sheaves}, Grund. Math.
Wiss. \textbf{265}, Springer-Verlag, Berlin, 1984.

\bibitem[GH1]{bib-GH1}
P. Griffiths and J. Harris, \emph{Principles of Algebraic Geometry},
Wiley-Interscience, New York, 1978.

\bibitem[GH2]{bib-GH2}
\sameauthor, ``On the Noether-Lefschetz theorem and some remarks on
codimension two cycles,'' Math. Ann. \textbf{271} (1985), 31--51.

\bibitem[Gr1]{bib-Gr1}
A. Grothendieck, ``Sur quelques points d'algèbre homologique,'' Tôhoku Math.
J. (2) \textbf{9} (1957), 119--221.

\bibitem[Gr2]{bib-Gr2}
\sameauthor, ``Sur la classification des fibrés holomorphes sur la sphère de
Riemann,'' Amer. J. Math. \textbf{79} (1957), 121--138.

\bibitem[Gr3]{bib-Gr3}
\sameauthor, ``The cohomology theory of abstract algebraic varieties,'' in
\emph{Proc. Internat. Congress Math. (Edinburgh, 1958)}, Cambridge U. P., New
York, 1960.

\bibitem[Gr4]{bib-Gr4}
\sameauthor, ``Fondements de la géométrie algébrique,'' in \emph{Collected
Bourbaki Talks}, Secrétariat Mathématique, Paris, 1962.

\bibitem[Gr5]{bib-Gr5}
\sameauthor, \emph{Récoltes et Semailles}, unpublished.

\bibitem[Gr-EGA]{bib-GrEGA}
\sameauthor{} (w. J. Dieudonné), \emph{Éléments de Géométrie Algébrique},
Publ. Math. IHES, Part I: \textbf{4} (1960); Part II: \textbf{8} (1961);
Part III: \textbf{11} (1961), \textbf{17} (1963); Part IV: \textbf{20} (1964),
\textbf{24} (1965), \textbf{28} (1966), \textbf{32} (1967).

\bibitem[Gr-EGA’]{bib-GrEGAI}
\sameauthor{} (w. J. Dieudonné), \emph{Éléments de Géométrie Algébrique I: Le
langage des schémas}, Grund. Math. Wiss. \textbf{166} (2nd ed.),
Springer-Verlag, Berlin-New York, 1971.

\bibitem[GrS]{bib-GrS}
A. Grothendieck and J.-P. Serre, \emph{Grothendieck-Serre Correspondence},
P. Colmez and J.-P. Serre ed., C. Maclean trans., Amer. Math. Soc. and Soc.
Math. de France, Paris, 2004.

\bibitem[H1]{bib-H1}
J. Harris, ``Galois groups of enumerative problems,'' Duke Math. J.
\textbf{46} (1979), no. \textbf{4}, 685--724.

\bibitem[H2]{bib-H2}
\sameauthor, \emph{Algebraic Geometry: A First Course}, Grad. Texts in Math.
\textbf{133}, Springer-Verlag, New York, 1995.

\bibitem[Ha1]{bib-Ha1}
R. Hartshorne, \emph{Algebraic Geometry}, Grad. Texts in Math. \textbf{52},
Springer-Verlag, New York-Heidelberg, 1977.

\bibitem[Ha2]{bib-Ha2}
\sameauthor, \emph{Ample Subvarieties of Algebraic Varieties}, Lect. Notes in
Math. \textbf{156}, Springer-Verlag, Berlin-New York, 1970.

\bibitem[Hat]{bib-Hat}
A. Hatcher, \emph{Algebraic Topology}, Cambridge U. P., Cambridge, 2002.

\bibitem[HM]{bib-HM}
M. Hazewinkel and C. Martin, ``A short elementary proof of Grothendieck's
theorem on algebraic vector bundles over the projective line,'' J. Pure Appl.
Algebra \textbf{25} (1982), no. \textbf{2}, 207--211.

\bibitem[He]{bib-He}
R. C. Heitmann, ``Examples of noncatenary rings,'' Trans. Amer. Math. Soc.
\textbf{247} (1979), 125--136.

\bibitem[Hil]{bib-Hil}
D. Hilbert, ``Ueber die Theorie der algebraischen Formen,'' Math. Ann.
\textbf{36} (1890), 473--534.

\bibitem[Hir]{bib-Hir}
H. Hironaka, ``Resolution of singularities of an algebraic variety over a
field of characteristic zero. I, II,'' Ann. of Math. (2) \textbf{79} (1964),
109-203; ibid. (2) \textbf{79} (1964), 205-326.

\bibitem[Ig]{bib-Ig}
J. Igusa, ``On some problems in abstract algebraic geometry,'' Proc. Nat.
Acad. Sci. \textbf{41} (1955), 964--967.

\bibitem[Ii]{bib-Ii}
S. Iitaka, \emph{Algebraic Geometry: An Introduction to Birational Geometry of
Algebraic Varieties}, Grad. Texts in Math. \textbf{76}, Springer-Verlag, New
York-Berlin, 1982.

\bibitem[Il]{bib-Il}
L. Illusie, ``Frobenius and Hodge degeneration,'' in \emph{Introduction to
Hodge Theory}, SMF/AMS Texts and Monographs \textbf{8}, 99--149, 2002.

\bibitem[dJ]{bib-dJ}
A. J. de Jong, ``Smoothness, semi-stability and alterations,'' IHES Publ.
Math. \textbf{83} (1996), 51--93.

\bibitem[KN]{bib-KN}
S. Kaji and M. Nakagawa, ``The Chow rings of the algebraic groups of type
$E_6$, $E_7$, and $E_8$,'' arXiv:0709.3702v3, preprint 2007.

\bibitem[KS1]{bib-KS1}
M. Kashiwara and P. Schapira, \emph{Sheaves on Manifolds}, Grund. Math. Wiss.
\textbf{292}, Springer-Verlag, Berlin, 1994.

\bibitem[KS2]{bib-KS2}
\sameauthor, \emph{Categories and sheaves}, Grund. Math. Wiss. \textbf{332},
Springer-Verlag, Berlin, 2006.

\bibitem[KPV]{bib-KPV}
K. Kedlaya, B. Poonen, and R. Vakil, \emph{The William Lowell Putnam
Mathematical Competition 1985--2000: Problems, Solutions, and Commentary},
Math. Ass. of Amer., Washington, DC, 2002.

\bibitem[Kl1]{bib-Kl1}
S. Kleiman, ``Toward a numerical theory of ampleness,'' Ann. Math. (2)
\textbf{84} (1966), no. \textbf{3}, 293--344.

\bibitem[Kl2]{bib-Kl2}
\sameauthor, ``The transversality of a general translate,'' Compositio Math.
\textbf{28} (1974), no. \textbf{3}, 287--297.

\bibitem[Kl3]{bib-Kl3}
\sameauthor, ``Cartier divisors versus invertible sheaves,'' Comm. Alg.
\textbf{28} (2000), no. \textbf{12}, 5677--5678.

\bibitem[Ko1]{bib-Ko1}
J. Kollár, \emph{Rational Curves on Algebraic Varieties}, Ergeb. Math. Grenz.
\textbf{32}, Springer-Verlag, Berlin, 1996.

\bibitem[Ko2]{bib-Ko2}
\sameauthor, \emph{Lectures on Resolutions of Singularities}, Annals of Math.
Stud. \textbf{166}, Princeton U. P., Princeton, NJ, 2007.

\bibitem[Lan]{bib-Lan}
S. Lang, \emph{Algebra} (rev. 3rd ed.), Grad. Texts in Math. \textbf{211},
Springer-Verlag, New York, 2002.

\bibitem[Laz]{bib-Laz}
R. Lazarsfeld, \emph{Positivity in Algebraic Geometry I: Classical Setting:
Line Bundles and Linear Series}, Ergeb. Math. Grenz. \textbf{48},
Springer-Verlag, Berlin, 2004.

\bibitem[Lef]{bib-Lef}
S. Lefschetz, \emph{L'analysis situs et la géométrie algébrique},
Gauthier-Villars, Paris, 1924.

\bibitem[Lew]{bib-Lew}
R. Lewontin, \emph{The Triple Helix: Gene, Organism, and Environment}, Harvard
U. P., Cambridge MA, 2000.

\bibitem[Lit]{bib-Lit}
J. E. Littlewood, \emph{Littlewood's Miscellany}, Cambridge U. P., Cambridge,
1986.

\bibitem[Liu]{bib-Liu}
Q. Liu, \emph{Algebraic Geometry and Arithmetic Curves}, R. Erné trans.,
Oxford Grad. Texts in Math. \textbf{6}, Oxford U. P., Oxford, 2002.

\bibitem[Lur]{bib-Lur}
J. Lurie, \emph{Higher Algebra}, preprint, September 18, 2017.

\bibitem[Lüt]{bib-Lüt}
W. Lütkebohmert, ``On compactification of schemes,'' Manuscripta Math.
\textbf{80} (1), 95--111.\footnote{原书此处未标年份，兹照录；通行著录为 1993
年。——译者注}

\bibitem[Mac]{bib-Mac}
S. Mac Lane, \emph{Categories for the Working Mathematician}, 2nd ed., Grad.
Texts in Math. \textbf{5}, Springer-Verlag, New York, 1998.

\bibitem[Mat1]{bib-Mat1}
H. Matsumura, \emph{Commutative Algebra}, 2nd ed., Math. Lecture Note Series
\textbf{56}, Benjamin/Cummings Publ. Co., Inc., Reading, MA, 1980.

\bibitem[Mat2]{bib-Mat2}
\sameauthor, \emph{Commutative Ring Theory}, M. Reid trans., 2nd ed.,\footnote
{原书此处作『2nd ed,』（句点脱漏），当为『2nd ed.,』之误，径改。——译者注}
Cambridge Stud. in Adv. Math. \textbf{8}, Cambridge U. P., Cambridge, 1989.

\bibitem[MO33489]{bib-MO33489}
\url{http://mathoverflow.net/questions/33489/}

\bibitem[MO68421]{bib-MO68421}
Answers by G. Elencwajg and M. Kim to
\url{http://mathoverflow.net/questions/68421/}

\bibitem[MO70920]{bib-MO70920}
Answers by T. Ekedahl to \url{http://mathoverflow.net/questions/70920/}

\bibitem[MO90551]{bib-MO90551}
Answer by M. Brandenburg to
\url{http://mathoverflow.net/questions/90551/}

\bibitem[MO129242]{bib-MO129242}
Answer by L. Moret-Bailly to
\url{http://mathoverflow.net/questions/129242/}

\bibitem[MO152008]{bib-MO152008}
Answer by J. Stillwell to
\url{http://mathoverflow.net/questions/152008/}

\bibitem[MO286390]{bib-MO286390}
Answers to \url{https://mathoverflow.net/questions/286390/}

\bibitem[Mc]{bib-Mc}
C. McLarty, ``The Rising Sea: Grothendieck on simplicity and generality,'' in
J. J. Gray and K. H. Parshall eds., \emph{Episodes in the History of Modern
Algebra (1800--1950)}, Amer. Math. Soc., Providence, RI, 2007.

\bibitem[Mu1]{bib-Mu1}
D. Mumford, ``Pathologies of modular algebraic surfaces,'' Amer. J. Math.
\textbf{83} (no. \textbf{2}), Apr. 1961, 339--342.

\bibitem[Mu2]{bib-Mu2}
\sameauthor, \emph{Lectures on Curves on an Algebraic Surface}, Princeton U.
P., Princeton NJ, 1966.

\bibitem[Mu3]{bib-Mu3}
\sameauthor, \emph{Abelian Varieties}, TIFR Studies in Math. \textbf{5},
Oxford U. P., London, 1970.

\bibitem[Mu4]{bib-Mu4}
\sameauthor, \emph{Curves and Their Jacobians}, Univ. of Michigan Press, Ann
Arbor, MI, 1975.

\bibitem[Mu5]{bib-Mu5}
\sameauthor, ``Hilbert's fourteenth problem --- The finite generation of
subrings such as rings of invariants,'' in \emph{Mathematical Developments
Arising from Hilbert Problems, Part 2}, F. Browder ed., Amer. Math. Soc.,
Providence RI, 1976.

\bibitem[Mu6]{bib-Mu6}
\sameauthor, \emph{Algebraic Geometry I: Complex Projective Varieties},
Classics in Math., Springer-Verlag, Berlin, 1995.

\bibitem[Mu7]{bib-Mu7}
\sameauthor, \emph{The Red Book of Varieties and Schemes} (2nd expanded ed.),
Lect. Notes in Math. \textbf{1358}, Springer-Verlag, Berlin, 1999. (Caution:
In order to keep readers on their toes, a number of typos were introduced
between the first and second editions.)

\bibitem[Mur]{bib-Mur}
D. Murfet, \url{http://therisingsea.org/notes/ModulesOverAScheme.pdf}

\bibitem[N]{bib-N}
F. Nietzsche, \emph{Twilight of the Idols, or, How to Philosophize with a
Hammer}, 1889.

\bibitem[OB]{bib-OB}
A. Ogus and G. Bergman, ``Nakayama's lemma for half-exact functors,'' Proc.
Amer. Math. Soc. \textbf{31} (1972), 67--74.

\bibitem[OSS]{bib-OSS}
C. Okonek, M. Schneider, and H. Spindler, \emph{Vector Bundles on Complex
Projective Spaces}, corrected reprint of the 1988 edition,
Birkhäuser/Springer Basel AG, Basel, 2011.

\bibitem[Pi]{bib-Pi}
P. Picasso, ``Statement to Marius de Zayas,'' 1923, in \emph{Picasso Speaks},
The Arts, New York, May 1923, pp. 315-26.

\bibitem[Po1]{bib-Po1}
H. Poincaré, \emph{Analysis situs}, J. de l'École Polytechnique (2)
\textbf{1} (1895), 1--123.

\bibitem[Po2]{bib-Po2}
\sameauthor, \emph{Papers on topology: Analysis situs and its five
supplements}, J. Stillwell trans., Amer. Math. Soc., Providence, RI, 2010.

\bibitem[P1]{bib-P1}
B. Poonen, ``Varieties without extra automorphisms, I: Curves,'' Math. Res.
Lett. \textbf{7} (2000), no. \textbf{1}, 67--76.

\bibitem[P2]{bib-P2}
\sameauthor, \emph{Rational Points on Varieties}, Grad. Stud. Math.,
\textbf{186}, Amer. Math. Soc., Providence, RI, 2017.

\bibitem[Ra]{bib-Ra}
M. Raynaud, ``Contre-exemple au vanishing theorem en caractéristique $p>0$,''
in \emph{C. P. Ramanujam --- A Tribute}, 273--278, TIFR Studies in Math.
\textbf{8}, Springer, Berlin-New York, 1978.

\bibitem[RS]{bib-RS}
M. Reed and B. Simon, \emph{Methods of Modern Mathematical Physics I:
Functional Analysis}, Academic Press, Inc., New York, 1980.

\bibitem[Rei1]{bib-Rei1}
M. Reid, \emph{Undergraduate Algebraic Geometry}, London Math. Soc. Student
Texts \textbf{12}, Cambridge U. P., Cambridge, 1988.

\bibitem[Rei2]{bib-Rei2}
\sameauthor, \emph{Undergraduate Commutative Algebra}, London Math. Soc.
Student Texts \textbf{29}, Cambridge U. P., Cambridge, 1995.

\bibitem[Rem]{bib-Rem}
R. Remmert, ``Local theory of complex spaces,'' in \emph{Several Complex
Variables VII: Sheaf-Theoretical Methods in Complex Analysis}, Encl. of Math.
Sci. \textbf{74}, Springer-Verlag, Berlin, 1994.

\bibitem[Roh]{bib-Roh}
F. Rohrer, ``Quasicoherent sheaves on toric schemes,'' Expo. Math.
\textbf{32} (2014), 33--78.

\bibitem[RotV]{bib-RotV}
M. Roth and R. Vakil, ``The affine stratification number and the moduli space
of curves,'' CRM Proceedings and Lecture Notes \textbf{38}, Univ. de
Montréal, 2004, 213--227.

\bibitem[Sa]{bib-Sa}
P. Salmon, ``Sulla fattorialità delle algebre graduate e degli anelli
locali,'' Rend. Sem. Mat. Univ. Padova \textbf{41} (1968), 119--138.

\bibitem[Schm]{bib-Schm}
H. L. Schmid, ``Über die Automorphismen eines algebraischen
Funktionenkörpers von Primzahlcharakteristik,'\footnote{原书此处作
『Primzahlcharacteristik』，当为『Primzahlcharakteristik』之误，径改。——译者注}
J. Reine Angew. Math. \textbf{179} (1938), 5--14.

\bibitem[Se1]{bib-Se1}
J.-P. Serre, ``Faisceaux algébriques cohérents,'' Ann. of Math. (2)
\textbf{61} (1955), 197--278.

\bibitem[Se2]{bib-Se2}
\sameauthor, ``Sur la dimension cohomologique des anneaux et des modules
noethériens,'' in \emph{Proc. Intern. Symp. on Alg. Number Theory}, 176--189,
Tokyo-Nikko, 1955.

\bibitem[Se3]{bib-Se3}
\sameauthor, ``Géométrie algébrique et géométrie analytique,'' Ann. de
l'institut Fourier \textbf{6} (1956), 1--42.

\bibitem[Se4]{bib-Se4}
\sameauthor, ``Sur la topologie des variétés algébriques en caractéristique
$p$,'' in \emph{Symposium Internacional de Topología Algebraica},
International Symposium on Algebraic Topology, 24--53, Univ. Nacional
Autónoma de México and UNESCO, Mexico City, 1958.

\bibitem[Se5]{bib-Se5}
\sameauthor, \emph{A Course in Arithmetic}, Grad. Texts in Math. \textbf{7},
Springer-Verlag, New York, 1973.

\bibitem[SGA2]{bib-SGA2}
A. Grothendieck, \emph{Séminaire de Géométrie Algébrique du Bois Marie 1962:
Cohomologie locale des faisceaux cohérents et Théorèmes de Lefschetz locaux et
globaux} (SGA2), North-Holland Publ. Co., Amsterdam, 1968.

\bibitem[SGA6]{bib-SGA6}
P. Berthelot, A. Grothendieck, and L. Illusie eds., \emph{Séminaire de
Géométrie Algébrique du Bois Marie 1966--67: Théorie des intersections et
théorème de Riemann Roch} (SGA 6), Lect. Notes in Math. \textbf{225},
Springer-Verlag, Berlin-New York, 1971.

\bibitem[Sh]{bib-Sh}
I. Shafarevich, \emph{Basic Algebraic Geometry 1: Varieties in Projective
Space}, 2nd ed., M. Reid trans., Springer-Verlag, Berlin, 1994.

\bibitem[Si1]{bib-Si1}
J. Silverman, \emph{The Arithmetic of Elliptic Curves}, 2nd ed., Grad. Texts
in Math. \textbf{206}, Springer-Verlag, New York-Berlin, 2009.

\bibitem[Si2]{bib-Si2}
\sameauthor, \emph{Advanced Topics in the Arithmetic of Elliptic Curves},
Grad. Texts in Math. \textbf{151}, Springer-Verlag, New York-Berlin, 1994.

\bibitem[Stacks]{bib-Stacks}
The Stacks Project Authors, \emph{Stacks Project},
\url{http://stacks.math.columbia.edu}. (To look up tags:
\url{http://stacks.math.columbia.edu/tags}. To search:
\url{http://stacks.math.columbia.edu/search}.)

\bibitem[Stic]{bib-Stic}
H. Stichtenoth, ``Über die Automorphismengruppe eines algebraischen
Funktionenkörpers von Primzahlcharakteristik. I. Eine Abschätzung der Ordnung
der Automorphismengruppe,'' Arch. Math. (Basel) \textbf{24} (1973), 527--544.

\bibitem[Ta]{bib-Ta}
A. Tarski, \emph{A decision method for elementary algebra and geometry}, RAND
Corp., Santa Monica, CA, 1948.

\bibitem[TWi]{bib-TWi}
R. Taylor and A. Wiles, ``Ring-theoretic properties of certain Hecke
algebras,'' Ann. of Math. (2) \textbf{141} (\textbf{3}) (1995), 553--572.

\bibitem[To]{bib-To}
L. Tolstoy, \emph{Anna Karenina},
\url{http://www.gutenberg.org/files/1399/1399-h/1399-h.htm}.

\bibitem[Tu]{bib-Tu}
L. Tu, \emph{An Introduction to Manifolds}, Universitext, Springer-Verlag,
New York, 2011.

\bibitem[Vak1]{bib-Vak1}
R. Vakil, \emph{The Rising Sea: Foundations of Algebraic Geometry}, preprint.

\bibitem[Vak2]{bib-Vak2}
\sameauthor, \emph{Puzzling through exact sequences: A bedtime story with
pictures},
\url{https://www.3blue1brown.com/blog/exact-sequence-picturebook}.

\bibitem[van]{bib-van}
B. L. van der Waerden, \emph{Algebra} vol. II, J. R. Schulenberger trans.,
Springer-Verlag, New York, 1991.

\bibitem[Var]{bib-Var}
V. S. Varadarajan, \emph{Lie Groups, Lie Algebras, and Their Representations},
Grad. Texts in Math. \textbf{102}, Springer-Verlag, New York-Berlin, 1984.

\bibitem[Vo]{bib-Vo}
C. Voisin, \emph{Hodge Theory and Complex Algebraic Geometry I and II},
L. Schneps trans., Cambridge Stud. in Adv. Math. \textbf{76} and
\textbf{77}, Cambridge U. P., Cambridge, 2007.

\bibitem[Weib]{bib-Weib}
C. Weibel, \emph{An Introduction to Homological Algebra}, Cambridge Stud. in
Adv. Math. \textbf{38}, Cambridge U. P., Cambridge, 1994.

\bibitem[Weil]{bib-Weil}
A. Weil, \emph{Scientific Works --- Collected Papers Vol. II} (1951--1964),
Springer-Verlag, New York-Heidelberg, 1979.

\bibitem[Wey]{bib-Wey}
H. Weyl, \emph{Philosophy of Mathematics and Natural Science}, Princeton U.
P., Princeton, NJ, 2009.

\bibitem[Wi]{bib-Wi}
A. Wiles, ``Modular elliptic curves and Fermat's last theorem,'' Ann. of
Math. (2) \textbf{141} (\textbf{3}) (1995), 443--551.

\bibitem[Wł]{bib-Wlodarczyk}
J. Włodarczyk, ``Toroidal varieties and the weak factorization theorem,''
Invent. Math. \textbf{154} (2003), no. \textbf{2}, 223--331.

\bibitem[Z]{bib-Z}
O. Zariski, ``Foundations of a general theory of birational
correspondences,'' Trans. Amer. Math. Soc. \textbf{53} (1943), 490--542.

\end{thebibliography}
\endgroup
